\documentclass[11pt]{article}
\usepackage[margin=1in]{geometry}
\usepackage{amsmath, amssymb, amsthm}
\usepackage{enumitem}
\usepackage{titlesec}
\usepackage{fancyhdr}
\usepackage{xcolor}
\usepackage{tcolorbox}
\usepackage{booktabs}
\usepackage{hyperref}

% Colors
\definecolor{creation}{RGB}{34, 85, 136}
\definecolor{concept}{RGB}{140, 60, 60}
\definecolor{practice}{RGB}{50, 110, 50}

% Box environments
\newtcolorbox{creationbox}[1][]{
  colback=creation!5, colframe=creation!80, fonttitle=\bfseries,
  title={#1}, arc=2mm
}

\newtcolorbox{conceptbox}[1][]{
  colback=concept!5, colframe=concept!70, fonttitle=\bfseries,
  title={#1}, arc=2mm
}

\newtcolorbox{practicebox}[1][]{
  colback=practice!5, colframe=practice!70, fonttitle=\bfseries,
  title={#1}, arc=2mm
}

\newtheorem{theorem}{Theorem}
\newtheorem{definition}{Definition}

\pagestyle{fancy}
\fancyhf{}
\rhead{ORF309/MAT380}
\lhead{Bernoulli Processes: A Creation Story}
\cfoot{\thepage}

\titleformat{\section}{\Large\bfseries\color{creation}}{}{0em}{}
\titleformat{\subsection}{\large\bfseries\color{concept}}{}{0em}{}

\begin{document}

\begin{center}
{\LARGE\bfseries\color{creation} In the Beginning, There Was Randomness}\\[6pt]
{\large A Creation Story for Bernoulli Processes II \& III}\\[4pt]
{\normalsize ORF309/MAT380 --- Lectures 08 \& 09 Study Guide}
\end{center}

\vspace{0.5cm}

\begin{creationbox}[The Premise]
In Genesis, creation follows a pattern: \textbf{Days 1--3} form empty domains (light/dark, sky/sea, land/vegetation), and \textbf{Days 4--6} fill those domains with agents (luminaries, birds/fish, animals/humans). Day 7 is rest---completion.

We use this structure to build up the probability concepts from Lectures 08--09. Each ``day'' introduces a foundational idea, and later days fill those foundations with power and meaning.
\end{creationbox}

\vspace{0.5cm}

%==============================================================================
\section{Day 1 --- Light from Darkness: The Bernoulli Trial}
%==============================================================================

\begin{creationbox}[Symbolic Parallel]
Day 1 separates light from darkness---the most basic binary distinction. A Bernoulli trial is the same act: every experiment has exactly two outcomes, success or failure, separated cleanly.
\end{creationbox}

\textbf{Recall.} A Bernoulli random variable $X_i$ takes values in $\{0, 1\}$:
\[
P(X_i = 1) = p, \qquad P(X_i = 0) = q = 1-p
\]

A \textbf{Bernoulli Process} sums $n$ independent trials:
\[
S_n = X_1 + X_2 + \cdots + X_n \sim \text{Binom}(n, p)
\]
\[
P(S_n = k) = \binom{n}{k} p^k q^{n-k}
\]

\begin{conceptbox}[Key Intuition]
$S_n$ counts \emph{how many} successes occurred in $n$ trials. It answers a ``how many?'' question. But it says nothing about \emph{when} those successes arrived. That's what comes next.
\end{conceptbox}

%==============================================================================
\section{Day 2 --- The Firmament: Arrival Times ($T_k$)}
%==============================================================================

\begin{creationbox}[Symbolic Parallel]
Day 2 creates structure---the firmament that separates waters above from waters below. Here we create the structural framework for \emph{when} events happen: arrival times. Like the firmament, this is architecture before content.
\end{creationbox}

\begin{definition}[Time of the $k$-th Success]
$T_k$ = the trial number on which the $k$-th success occurs.
\end{definition}

\textbf{Example timeline:}
\begin{center}
\begin{tabular}{ccccccccccc}
f & f & \textbf{s} & f & f & f & f & \textbf{s} & \textbf{s} & f & \textbf{s} \\
1 & 2 & 3 & 4 & 5 & 6 & 7 & 8 & 9 & 10 & 11
\end{tabular}
\end{center}
Here $T_1 = 3$, $T_2 = 8$, $T_3 = 9$, $T_4 = 11$.

\subsection{Deriving $P(T_k = n)$}

The event $\{T_k = n\}$ means: exactly $k-1$ successes in the first $n-1$ trials, \emph{and} a success on trial $n$. These are independent, so:

\[
\boxed{P(T_k = n) = \underbrace{P(S_{n-1} = k-1)}_{\text{$k-1$ successes in $n-1$ trials}} \cdot \underbrace{P(X_n = 1)}_{\text{success on trial } n} = \binom{n-1}{k-1} p^k \, q^{n-k}}
\]

for $n \geq k$.

\begin{definition}[Negative Binomial Distribution]
A DRV $Z$ with $P(Z = n) = \binom{n-1}{k-1} p^k q^{n-k}$ for $n = k, k+1, k+2, \ldots$ is said to have the \textbf{Negative Binomial distribution}. Notation: $Z \sim \text{NegBinom}(k, p)$.
\end{definition}

%==============================================================================
\section{Day 3 --- Vegetation and Self-Reproduction: The Geometric Distribution}
%==============================================================================

\begin{creationbox}[Symbolic Parallel]
Day 3 creates life that reproduces itself---vegetation ``bearing seed according to its kind.'' The Geometric distribution is the most elemental self-reproducing unit of the Bernoulli process: the time to the \emph{first} success. Every subsequent inter-arrival time is a fresh copy of this same distribution. Like a seed, once you have the Geometric, the entire process regenerates from it.
\end{creationbox}

\begin{definition}[Geometric Distribution]
The special case $k = 1$ of the Negative Binomial. If $T_1$ = time until the first success:
\[
\boxed{P(T_1 = n) = p \, q^{n-1}, \qquad n = 1, 2, 3, \ldots}
\]
Notation: $T_1 \sim \text{Geom}(p)$.
\end{definition}

\subsection{Expectation of $T_1$}

\[
E[T_1] = \sum_{k=1}^{\infty} k \cdot p q^{k-1} = p \sum_{k=1}^{\infty} k \, q^{k-1} = p \cdot \frac{d}{dq}\left(\sum_{k=1}^{\infty} q^k\right) = p \cdot \frac{d}{dq}\left(\frac{1}{1-q}\right) = p \cdot \frac{1}{(1-q)^2} = \frac{1}{p}
\]

\begin{conceptbox}[Key Result]
\[
\boxed{E[T_1] = \frac{1}{p}}
\]
If $p = 0.01$, expect to wait 100 trials for your first success. The rarer the event, the longer the wait.
\end{conceptbox}

\subsection{Can You Fail Forever?}

If $p > 0$:
\[
P(T_1 = \infty) = P(X_1 = 0, X_2 = 0, X_3 = 0, \ldots) = q \cdot q \cdot q \cdots = q^\infty = 0
\]

Therefore $P(T_1 < \infty) = 1$. \textbf{If there is any chance of success, you will eventually succeed.}

\begin{conceptbox}[Murphy's Law (Probabilistic Version)]
Flip the framing: if $p$ is the probability of \emph{failure}, then $P(\text{failure eventually happens}) = 1$. Whatever can go wrong, will go wrong.

\textbf{Bonus:} This proves the geometric series without calculus:
\[
\sum_{k=1}^{\infty} p q^{k-1} = 1 \implies \sum_{k=0}^{\infty} q^k = \frac{1}{1-q}
\]
\end{conceptbox}

%==============================================================================
\section{Day 4 --- Luminaries Governing Time: Inter-Arrival Times}
%==============================================================================

\begin{creationbox}[Symbolic Parallel]
Day 4 places sun, moon, and stars to \emph{govern} the light/dark cycle of Day 1---giving time its rulers and rhythm. Here we discover that the gaps between successes are governed by the same Geometric distribution. The ``clock resets'' after each success, giving the process its temporal rhythm.
\end{creationbox}

\subsection{The Inter-Arrival Structure}

Define the gaps between consecutive successes:
\[
T_1, \quad T_2 - T_1, \quad T_3 - T_2, \quad \ldots
\]

\textbf{Each gap is independently $\text{Geom}(p)$.} The process has no memory---after each success, it's as if you're starting fresh.

\textbf{Joint probability example:}
\begin{align*}
&P(T_1 = 3,\; T_2 - T_1 = 5,\; T_3 - T_2 = 1) \\
&= P(\text{ffs, ffffs, s}) = (q^2 p)(q^4 p)(p) \\
&= P(T_1 = 3) \cdot P(T_1 = 5) \cdot P(T_1 = 1)
\end{align*}

The joint distribution \emph{factors} into a product of Geometric distributions.

\subsection{Telescopic Sum and $E[T_k]$}

Since $T_k = T_1 + (T_2 - T_1) + (T_3 - T_2) + \cdots + (T_k - T_{k-1})$, and each term is independently $\text{Geom}(p)$:

\[
\boxed{E[T_k] = k \cdot E[T_1] = \frac{k}{p}}
\]

\begin{conceptbox}[Key Result]
The expected time to the $k$-th success is $k/p$. To get 5 successes when $p = 0.1$, expect to wait $5/0.1 = 50$ trials on average.
\end{conceptbox}

\subsection{Can You Succeed Infinitely Often?}

If $p > 0$:
\[
P(\text{success occurs infinitely often}) = P(T_1 < \infty) \cdot P(T_2 - T_1 < \infty) \cdot P(T_3 - T_2 < \infty) \cdots = 1
\]

\textbf{Strong Murphy's Law:} If something can happen, it will happen infinitely many times.

%==============================================================================
\section{Day 5 --- Teeming Life: The Superexperiment}
%==============================================================================

\begin{creationbox}[Symbolic Parallel]
Day 5 fills the empty sky and seas with teeming life---birds and fish, multiplying abundantly. The Superexperiment does the same: we repeat the original experiment not once or twice, but \emph{many} times, filling the sample space with data. The question shifts from individual trials to the behavior of crowds.
\end{creationbox}

\subsection{Setup}

Take an original experiment with event $A$ where $P(A) = p$. Repeat it $n$ times independently. Define:
\[
S_n = \sum_{k=1}^{n} \mathbf{1}_{A_k} \sim \text{Binom}(n, p)
\]

The fraction of times $A$ occurred:
\[
\frac{S_n}{n} = \frac{1}{n} \sum_{k=1}^{n} \mathbf{1}_{A_k}
\]

This is a \textbf{random variable} taking values in $[0, 1]$.

\begin{conceptbox}[Key Property]
The expectation of $S_n/n$ is constant for all $n$:
\[
E\!\left[\frac{S_n}{n}\right] = \frac{1}{n} \sum_{k=1}^{n} P(A_k) = \frac{1}{n} \cdot np = p
\]
The average fraction always equals $p$, regardless of how many trials you run.
\end{conceptbox}

As $n$ grows, the histogram of $S_n/n$ concentrates tighter and tighter around $p$. The randomness \emph{shrinks}.

%==============================================================================
\section{Day 6 --- Humanity and Dominion: The Law of Large Numbers}
%==============================================================================

\begin{creationbox}[Symbolic Parallel]
Day 6 creates humanity in the image of God and gives them dominion---the power to name, order, and steward creation. The Law of Large Numbers is the ``dominion'' theorem of probability: it tells us that \emph{in the limit}, randomness submits to order. The random fraction $S_n/n$ converges to the deterministic value $p$. Chaos becomes knowable.
\end{creationbox}

\begin{theorem}[Law of Large Numbers]
With probability 1:
\[
\boxed{\lim_{n \to \infty} \frac{S_n}{n} = \lim_{n \to \infty} \frac{1}{n} \sum_{k=1}^{n} \mathbf{1}_{A_k} = P(A) = p}
\]
\end{theorem}

\textbf{What this means:}
\begin{itemize}[nosep]
  \item For any finite $n$: $S_n/n$ is \textbf{random}.
  \item In the limit $n \to \infty$: $\lim S_n/n$ is \textbf{not random}---it equals exactly $p$.
  \item Probabilities \emph{are} long-run frequencies, but \textbf{only} at infinity.
\end{itemize}

\subsection{Proof via Variance}

To prove LLN, we need to show the randomness of $S_n/n$ vanishes. This requires \textbf{variance}.

\begin{definition}[Variance]
For a random variable $X$:
\[
\text{Var}(X) = E\!\left[(X - E[X])^2\right]
\]
Variance measures how far $X$ deviates from its expectation. If $\text{Var}(X) = 0$, then $X$ is not random.
\end{definition}

\textbf{Properties of Variance:}
\begin{enumerate}[nosep]
  \item $\text{Var}(X) \geq 0$
  \item $\text{Var}(X) = 0 \iff X$ is deterministic (not random)
  \item $\text{Var}(X) = E[X^2] - (E[X])^2$ \hfill (computational formula)
  \item If $X \perp Y$: $\text{Var}(X + Y) = \text{Var}(X) + \text{Var}(Y)$ \hfill (additivity for independent RVs)
  \item $\text{Var}(\alpha X) = \alpha^2 \, \text{Var}(X)$ \hfill (scaling)
\end{enumerate}

\subsection{Variance of an Indicator}

\[
\text{Var}(\mathbf{1}_{A_i}) = E[\mathbf{1}_{A_i}^2] - (E[\mathbf{1}_{A_i}])^2 = p - p^2 = p(1-p)
\]

\subsection{The Proof}

\begin{align*}
\text{Var}\!\left(\frac{S_n}{n}\right)
&= \text{Var}\!\left(\frac{1}{n}\sum_{k=1}^{n} \mathbf{1}_{A_k}\right)
= \frac{1}{n^2} \sum_{k=1}^{n} \text{Var}(\mathbf{1}_{A_k})  \tag{Properties 5 \& 4}\\
&= \frac{1}{n^2} \cdot n \cdot p(1-p) = \frac{p(1-p)}{n}
\end{align*}

Taking the limit:
\[
\lim_{n \to \infty} \text{Var}\!\left(\frac{S_n}{n}\right) = \lim_{n \to \infty} \frac{p(1-p)}{n} = 0
\]

Since the variance goes to zero, the randomness vanishes, and $S_n/n \to p$. $\qed$

\subsection{LLN for General Random Variables}

\begin{theorem}
For i.i.d.\ random variables $X_1, X_2, \ldots$ with common distribution $X$:
\[
\boxed{\lim_{n \to \infty} \frac{1}{n} \sum_{k=1}^{n} X_k = E[X] \quad \text{with probability 1}}
\]
\end{theorem}

Expectations are averages---but \textbf{only} when $n \to \infty$.

\subsection{LLN for Conditional Probability}

\[
\lim_{n \to \infty} \frac{\sum_{k=1}^{n} \mathbf{1}_{A_k \cap B_k}}{\sum_{k=1}^{n} \mathbf{1}_{B_k}} = \frac{P(A \cap B)}{P(B)} = P(A \mid B)
\]

%==============================================================================
\section{Day 7 --- Rest: The Completed Framework}
%==============================================================================

\begin{creationbox}[Symbolic Parallel]
Day 7 is rest---not inactivity, but the recognition that the structure is complete. Here we step back and see the whole architecture. Everything connects.
\end{creationbox}

\renewcommand{\arraystretch}{1.4}
\begin{center}
\begin{tabular}{@{}lll@{}}
\toprule
\textbf{Concept} & \textbf{Distribution} & \textbf{Key Formula} \\
\midrule
\# successes in $n$ trials & $\text{Binom}(n,p)$ & $P(S_n = k) = \binom{n}{k}p^k q^{n-k}$ \\[3pt]
Trial of $k$-th success & $\text{NegBinom}(k,p)$ & $P(T_k = n) = \binom{n-1}{k-1}p^k q^{n-k}$ \\[3pt]
Trial of 1st success & $\text{Geom}(p)$ & $P(T_1 = n) = pq^{n-1}$ \\[3pt]
Expected 1st success & --- & $E[T_1] = 1/p$ \\[3pt]
Expected $k$-th success & --- & $E[T_k] = k/p$ \\[3pt]
Variance of indicator & --- & $\text{Var}(\mathbf{1}_A) = p(1-p)$ \\[3pt]
LLN & --- & $\lim S_n/n = p$ \\
\bottomrule
\end{tabular}
\end{center}

\subsection{Sampling Methods Reference}

\renewcommand{\arraystretch}{1.3}
\begin{center}
\begin{tabular}{@{}lcc@{}}
\toprule
& \textbf{Ordered} & \textbf{Unordered} \\
\midrule
\textbf{Without replacement} & $\dfrac{n!}{(n-r)!}$ & $\dbinom{n}{r}$ \\[12pt]
\textbf{With replacement} & $n^r$ & $\dbinom{n+r-1}{r}$ \\[8pt]
\bottomrule
\end{tabular}
\end{center}

%==============================================================================
\section{Pset Preparation: Practice Problems}
%==============================================================================

\begin{practicebox}[Problem 1: Geometric Distribution]
You apply to jobs, and each application independently has a 5\% chance of resulting in an offer.
\begin{enumerate}[label=(\alph*), nosep]
  \item What is the probability your first offer comes on exactly the 10th application?
  \item What is the expected number of applications until your first offer?
  \item What is the probability you get \emph{no} offers in your first 20 applications?
\end{enumerate}
\end{practicebox}

\textbf{Approach:}
\begin{itemize}[nosep]
  \item Let $T_1$ = application number of first offer. $T_1 \sim \text{Geom}(0.05)$.
  \item (a): $P(T_1 = 10) = (0.05)(0.95)^{9}$
  \item (b): $E[T_1] = 1/0.05 = 20$
  \item (c): $P(T_1 > 20) = P(\text{20 failures}) = (0.95)^{20}$
\end{itemize}

\vspace{0.4cm}

\begin{practicebox}[Problem 2: Negative Binomial]
A basketball player makes free throws with probability $p = 0.7$.
\begin{enumerate}[label=(\alph*), nosep]
  \item What is the probability that her 3rd made free throw occurs on her 5th attempt?
  \item What is the expected number of attempts to make 3 free throws?
\end{enumerate}
\end{practicebox}

\textbf{Approach:}
\begin{itemize}[nosep]
  \item $T_3 \sim \text{NegBinom}(3, 0.7)$.
  \item (a): $P(T_3 = 5) = \binom{4}{2}(0.7)^3(0.3)^2$
  \item (b): $E[T_3] = 3/0.7 \approx 4.29$
\end{itemize}

\vspace{0.4cm}

\begin{practicebox}[Problem 3: Binomial and Complement Events]
A factory produces items with a 2\% defect rate, independently. In a batch of 100 items:
\begin{enumerate}[label=(\alph*), nosep]
  \item What is the probability of \emph{zero} defects?
  \item What is the probability of \emph{at least one} defect?
  \item What is the probability of \emph{at most 3} defects?
\end{enumerate}
\end{practicebox}

\textbf{Approach:}
\begin{itemize}[nosep]
  \item $S_{100} \sim \text{Binom}(100, 0.02)$.
  \item (a): $P(S_{100} = 0) = (0.98)^{100}$
  \item (b): $P(S_{100} \geq 1) = 1 - (0.98)^{100}$ \quad (complement!)
  \item (c): $P(S_{100} \leq 3) = \sum_{k=0}^{3} \binom{100}{k}(0.02)^k(0.98)^{100-k}$
\end{itemize}

\vspace{0.4cm}

\begin{practicebox}[Problem 4: LLN and Variance]
Let $X_1, X_2, \ldots$ be i.i.d.\ Bernoulli$(p)$ with $p = 0.3$.
\begin{enumerate}[label=(\alph*), nosep]
  \item Compute $\text{Var}(S_n/n)$ as a function of $n$.
  \item For what value of $n$ is $\text{Var}(S_n/n) < 0.001$?
  \item Explain in words what happens to $S_n/n$ as $n \to \infty$.
\end{enumerate}
\end{practicebox}

\textbf{Approach:}
\begin{itemize}[nosep]
  \item (a): $\text{Var}(S_n/n) = p(1-p)/n = 0.21/n$
  \item (b): Solve $0.21/n < 0.001 \implies n > 210$
  \item (c): By LLN, $S_n/n \to 0.3$ with probability 1. The sample average converges to the true probability.
\end{itemize}

\vspace{0.4cm}

\begin{practicebox}[Problem 5: VC-Style Problem (Comprehensive)]
A venture fund invests in startups sequentially. Each startup independently becomes a ``hit'' with probability $p = 0.08$.
\begin{enumerate}[label=(\alph*), nosep]
  \item What is the expected number of investments until the fund's first hit?
  \item The fund needs 4 hits to break even. What is the expected number of investments?
  \item What is the probability the fund gets its 4th hit on exactly the 50th investment?
  \item What is the probability the fund has at least 4 hits in its first 50 investments?
\end{enumerate}
\end{practicebox}

\textbf{Approach:}
\begin{itemize}[nosep]
  \item (a): $E[T_1] = 1/0.08 = 12.5$
  \item (b): $E[T_4] = 4/0.08 = 50$
  \item (c): $T_4 \sim \text{NegBinom}(4, 0.08)$. \; $P(T_4 = 50) = \binom{49}{3}(0.08)^4(0.92)^{46}$
  \item (d): Let $S_{50} \sim \text{Binom}(50, 0.08)$. \; $P(S_{50} \geq 4) = 1 - \sum_{k=0}^{3}\binom{50}{k}(0.08)^k(0.92)^{50-k}$
\end{itemize}

\vspace{0.4cm}

\begin{practicebox}[Problem 6: Murphy's Law and Infinite Successes]
Prove that if $p > 0$, success occurs infinitely often with probability 1.

\emph{Hint:} Write the event ``success occurs infinitely often'' in terms of $T_1, T_2 - T_1, T_3 - T_2, \ldots$ and use independence.
\end{practicebox}

\textbf{Approach:}
\begin{align*}
P(\text{success } \infty \text{ often}) &= P(T_1 < \infty, \; T_2 - T_1 < \infty, \; T_3 - T_2 < \infty, \ldots) \\
&= P(T_1 < \infty) \cdot P(T_1 < \infty) \cdot P(T_1 < \infty) \cdots \tag{independence, same distribution} \\
&= 1 \cdot 1 \cdot 1 \cdots = 1
\end{align*}

\vspace{0.4cm}

\begin{conceptbox}[Pset Strategy Checklist]
\begin{enumerate}[nosep]
  \item \textbf{Identify the random variable.} Is it counting successes ($S_n$, Binomial) or waiting for a success ($T_k$, Negative Binomial/Geometric)?
  \item \textbf{``At least one''} $\implies$ use the complement: $P(X \geq 1) = 1 - P(X = 0)$.
  \item \textbf{``At most $k$''} $\implies$ sum from 0 to $k$: $P(X \leq k) = \sum_{i=0}^{k} P(X = i)$.
  \item \textbf{Expected wait times} use telescopic sums: $E[T_k] = k/p$.
  \item \textbf{Variance questions} often use $\text{Var}(S_n/n) = p(1-p)/n$.
  \item \textbf{Independence} lets you multiply probabilities and add variances.
\end{enumerate}
\end{conceptbox}

\end{document}
